\begin{abstract}
Self-evolving agents convert execution feedback into persistent skills, but most studies treat the state-writing step as an implementation detail and focus instead on which trajectories should be shown to the updater. We ask whether this misses a distinct failure mode: the same evidence can be converted into materially different persistent states. We begin with a controlled DeepSeek-V4-Pro study of 48 matched evidence pairs, 96 learned states, and 1,728 held-out evaluations. Replacing winner evidence with a mixed rejected-witness view produces only a small heterogeneous mean gain of 0.023, with a frozen confidence interval that crosses zero. A single-stream mechanism study then separates evidence from state realization. A historically strong First-Fail state remains useful when frozen and re-evaluated, whereas two fresh updater realizations from reconstructed byte-identical evidence fail to reproduce that advantage. This combination localizes a candidate bottleneck at persistent-state generation rather than evidence availability alone, while leaving broad failure-evidence superiority unsupported. We formalize the pipeline as evidence selection $\rightarrow$ state generation $\rightarrow$ persistent state $\rightarrow$ downstream behavior and introduce a prospective variance-controlled typed compiler that maps learner-visible trajectory+score packages to a finite repair representation and canonical skill state. The planned bridge crosses \win{} versus First-Fail-4 evidence with native free-form versus constrained compilation on fresh development streams, including a score-matched trajectory-blind generic control. The present evidence supports a state-generation bottleneck hypothesis and a concrete causal test for it; it does not yet establish that the compiler improves utility or that rejected failures are uniquely beneficial.
\end{abstract}
